% LZ Field Cage — r-z half-plane cross-section
%
% Numbers hand-copied from data/lz_fc_barrels.csv,
% data/lz_fc_grids.csv, data/lxe_detector.csv as of 2026-05-03.
% If you change a CSV, ask Claude to resync this file (and its
% companion design/lz_fieldcage_model.md).
%
% Compile with:  pdflatex design/tikz/fieldcage.tex
% Output:        fieldcage.pdf  (one-page standalone diagram)

\documentclass[tikz,border=8pt]{standalone}
\usepackage{tikz}
\usetikzlibrary{calc,patterns}

% --- Dimensions in cm (single source of truth) -----------------------
% Barrel sources (data/lz_fc_barrels.csv)
\def\TpcR{72.8}        % FCPT inner radius = TPC inner
\def\RingR{74.3}       % FCRN/FCRS/FCSE inner radius = FC outer
\def\BarrelZbot{-13.75} % FC barrel z extent: bottom (RFR floor)
\def\BarrelZtop{145.6}  % FC barrel z extent: top (gate plane)

% Grid holders (data/lz_fc_grids.csv)
\def\HolderRin{72.8}   % annular inner = TPC inner
\def\HolderRout{80.3}  % annular outer = LXe-skin outer
\def\GateZ{145.6}      % FCTG face (slab opens upward)
\def\CathZ{0.0}        % FCBG face (slab opens downward)
\def\HfctgAxial{1.558} % FCTG slab axial thickness (back-derived)
\def\HfcbgAxial{0.779} % FCBG slab axial thickness (back-derived)

% Skin LXe outer (data/lxe_detector.csv)
\def\SkinOuter{82.1}   % R_ICV_inner

% FV box overlay (the bb0nu inner-volume default, docs/run_recipes.md)
\def\FvZmin{26.0}
\def\FvZmax{96.0}
\def\FvRmax{39.0}

\begin{document}
\begin{tikzpicture}[scale=0.07, every node/.style={font=\scriptsize}]

  % ==== Material / region styles ====================================
  \tikzset{
    lxe/.style       = {fill=cyan!10, draw=cyan!50!black, line width=0.2pt},
    skin/.style      = {fill=cyan!25, draw=cyan!50!black, line width=0.2pt},
    rfr/.style       = {fill=cyan!18, draw=cyan!50!black, line width=0.2pt},
    gas/.style       = {fill=yellow!8, draw=gray!60, line width=0.2pt},
    ti/.style        = {fill=gray!50, draw=black, line width=0.3pt},
    ss/.style        = {fill=gray!70, draw=black, line width=0.3pt},
    ptfe/.style      = {fill=brown!20, draw=brown!60!black, line width=0.3pt},
    ringline/.style  = {draw=gray!70!black, line width=0.4pt},
    fv/.style        = {draw=red!80!black, line width=0.6pt, dashed},
    dropped/.style   = {draw=gray!50, line width=0.4pt, dashed},
    dim/.style       = {<->, >=stealth, draw=black!70, line width=0.2pt},
  }

  % ==== Background LXe regions ======================================
  % Skin annulus: r ∈ [RingR, SkinOuter], z ∈ [BarrelZbot, GateZ]
  % (Note: Real geometry has skin starting at FCPT inner = TpcR; we draw the
  % portion outside the ring shell as cleanly-skin to avoid clutter.)
  \fill[skin] (\RingR, \BarrelZbot) rectangle (\SkinOuter, \GateZ);
  % Active TPC LXe: r ∈ [0, TpcR], z ∈ [CathZ, GateZ] (drift region)
  \fill[lxe]  (0, \CathZ) rectangle (\TpcR, \GateZ);
  % RFR (inert LXe, below cathode): r ∈ [0, TpcR], z ∈ [BarrelZbot, CathZ]
  \fill[rfr]  (0, \BarrelZbot) rectangle (\TpcR, \CathZ);
  % Gas above gate (inside the field cage but with no LXe interactions)
  \fill[gas]  (0, \GateZ) rectangle (\TpcR, \GateZ+18);

  % ==== FCPT — PTFE walls at R=72.8, full barrel ===================
  % Drawn at the realistic 5 mm thickness for visibility.
  \fill[ptfe] (\TpcR, \BarrelZbot) rectangle (\TpcR+0.5, \BarrelZtop);

  % ==== FC annulus (R=72.8 to 74.3) — transparent in MC ============
  % FCRN/FCRS/FCSE all sit at R=74.3 (FC outer); the 1.5 cm gap to PTFE
  % is the field-cage volume itself, treated as vacuum by the tracker.
  \draw[ringline] (\RingR, \BarrelZbot) -- (\RingR, \BarrelZtop);

  % ==== FCRN — Ti rings, drawn at the 65 individual ring positions ==
  % 25 mm pitch, with the ring shell smeared at R=74.3. Show 12 sample
  % rings as short tick marks for visual reference.
  \foreach \z in {-10, 0, 10, 20, 30, 40, 50, 60, 70, 80, 90, 100, 110, 120, 130, 140} {
    \draw[ti] (\RingR, \z-0.4) rectangle (\RingR+0.6, \z+0.4);
  }

  % ==== FCTG — top grids+holders annular slab =======================
  % Slab spans z ∈ [GateZ, GateZ + H], r ∈ [HolderRin, HolderRout]
  \fill[ss] (\HolderRin, \GateZ) rectangle (\HolderRout, \GateZ+\HfctgAxial);
  % LXe-facing face (bottom of slab) emphasised
  \draw[ringline, line width=0.6pt] (\HolderRin, \GateZ) -- (\HolderRout, \GateZ);

  % ==== FCBG — cathode grid+holder annular slab =====================
  % Slab spans z ∈ [-H, 0], r ∈ [HolderRin, HolderRout]
  \fill[ss] (\HolderRin, \CathZ-\HfcbgAxial) rectangle (\HolderRout, \CathZ);
  % LXe-facing face (top of slab) emphasised
  \draw[ringline, line width=0.6pt] (\HolderRin, \CathZ) -- (\HolderRout, \CathZ);

  % ==== Dropped bottom-shield grid (greyed out) =====================
  % Indicative location below RFR floor — drawn dashed/grey, with a
  % "dropped from model" label so the simplification is explicit.
  \draw[dropped] (\HolderRin, \BarrelZbot-2) rectangle (\HolderRout, \BarrelZbot-1);
  \node[anchor=west, gray!60!black] at (\HolderRout+1, \BarrelZbot-1.5)
    {\tiny dropped: bottom-shield grid};

  % ==== FV box overlay ==============================================
  \draw[fv] (0, \FvZmin) rectangle (\FvRmax, \FvZmax);
  \node[red!80!black, anchor=south west] at (1, \FvZmax-3) {FV};

  % ==== Axes ========================================================
  \draw[->, line width=0.3pt] (0, \BarrelZbot-15) -- (0, \GateZ+25)
    node[above] {$z$};
  \draw[->, line width=0.3pt] (0,0) -- (\SkinOuter+10, 0)
    node[right] {$r$};

  % ==== Annotations =================================================
  % Component callouts
  \node[anchor=west] at (\TpcR-10, 35)        {active LXe};
  \node[anchor=west] at (\TpcR+1.6, 60)       {\tiny FCPT (PTFE)};
  \node[anchor=west] at (\RingR+1.6, 75)      {\tiny FCRN/FCRS/FCSE};
  \node[anchor=west] at (\HolderRout+1, \GateZ+3)  {FCTG (anode+gate)};
  \node[anchor=west] at (\HolderRout+1, \CathZ-3)  {FCBG (cathode)};
  \node[anchor=west] at (\SkinOuter-15, 110)  {\tiny LXe skin};
  \node[anchor=west] at (5, \GateZ+10)        {gas};
  \node[anchor=west] at (\TpcR-15, -8)        {\tiny RFR (inert)};

  % z anchors on axis
  \node[anchor=east, fill=white, inner sep=1pt] at (-1, \GateZ)
    {gate $z=\GateZ$};
  \node[anchor=east, fill=white, inner sep=1pt] at (-1, \CathZ)
    {cathode $z=0$};
  \node[anchor=east, fill=white, inner sep=1pt] at (-1, \BarrelZbot)
    {RFR bot $z=\BarrelZbot$};

  % Radial dimensions
  \draw[dim] (0, \BarrelZbot-12) -- (\TpcR, \BarrelZbot-12);
  \node[below] at (\TpcR/2, \BarrelZbot-13) {$R_{\rm FC,inner} = \TpcR$ cm};

  \draw[dim] (0, \BarrelZbot-22) -- (\RingR, \BarrelZbot-22);
  \node[below] at (\RingR/2, \BarrelZbot-23)
    {$R_{\rm rings} = \RingR$ cm};

  \draw[dim] (\HolderRin, \GateZ+8) -- (\HolderRout, \GateZ+8);
  \node[above] at ((\HolderRin+\HolderRout)/2, \GateZ+8)
    {\tiny annulus $R \in [\HolderRin, \HolderRout]$};

\end{tikzpicture}
\end{document}
